% W5i87_NewFrontiers.tex
% DFD 20-Agent Research Programme --- Iteration 87
% New Frontiers, Paper Proposals, and Research Directions

\documentclass[11pt,a4paper]{article}
\usepackage[utf8]{inputenc}
\usepackage{amsmath,amssymb}
\usepackage{booktabs}
\usepackage{longtable}
\usepackage{geometry}
\usepackage{xcolor}
\usepackage{tcolorbox}
\usepackage{enumitem}
\geometry{margin=2.5cm}

\definecolor{dfdblue}{RGB}{0,50,180}

\title{W5i87: New Frontiers and Paper Proposals\\[6pt]
\large Three New Multi-Probe Channels $\cdot$ Textbook Pedagogical Strategy $\cdot$ 2702 Total Proposals}
\author{DFD 20-Agent Research Programme}
\date{2026-03-24}

\begin{document}
\maketitle

\section{Paper Proposals (i87 Batch: 110 new)}

\subsection{Tier 1: High Priority (4 proposals)}

\begin{enumerate}[nosep]
\item \textbf{``Baryonic Feedback on the DFD Signal: Implications for Euclid''} (10--12 pp, MNRAS). Quantifies the baryonic correction to DFD predictions across scales. Essential companion to the N-body paper.

\item \textbf{``Three New Probes of Density Field Dynamics: HOD Mismatch, Void Lensing, and kSZ''} (15 pp, JCAP). Develops the three new probes identified in Finding 228 into a quantitative multi-probe paper.

\item \textbf{``A Decision Tree for Euclid DR1: How to Interpret Results in the Context of DFD''} (8 pp, A\&A Letters). The formal pre-registration companion: specifies exactly what outcomes confirm, are inconclusive for, or falsify DFD.

\item \textbf{``DFD-CLASS: Design Document and Preliminary Benchmarks''} (CPC or Astron.\ Comp.). Scoping document for the modified Boltzmann code. Phase 1 implementation spec.
\end{enumerate}

\subsection{Tier 2: Medium Priority (15 proposals)}

\begin{enumerate}[nosep]
\setcounter{enumi}{4}
\item DFD and the ACT DR6 lensing power spectrum
\item DFD predictions for galaxy--galaxy lensing with Rubin/LSST
\item DFD and the thermal Sunyaev--Zel'dovich power spectrum
\item DFD void abundance function prediction
\item DFD and the ISW--void cross-correlation
\item DFD predictions for the dark energy spectroscopic instrument Y3
\item DFD and the galaxy angular momentum--mass relation
\item DFD predictions for cosmic shear peak counts (Euclid)
\item DFD and cluster gas fraction measurements
\item DFD predictions for the galaxy luminosity function evolution
\item DFD and the Lyman-$\alpha$ forest power spectrum
\item DFD predictions for photometric redshift distributions
\item DFD and the X-ray cluster temperature function
\item DFD predictions for the galaxy size--mass relation evolution
\item DFD and the parity-odd galaxy 4-point function
\end{enumerate}

\subsection{Tier 3: Exploratory (91 proposals)}

Topics include: DFD in Kaluza--Klein reduction, DFD and emergent spacetime, DFD and holographic entanglement entropy, DFD predictions for quantum clocks in gravitational fields, DFD and the Unruh effect, DFD implications for wormhole geometry, DFD and the information paradox, DFD predictions for post-merger gravitational wave spectra, DFD and neutron star equation of state, DFD implications for primordial magnetic fields, plus 81 additional exploratory directions.

\section{Multi-Probe Enhancement Analysis}

The three new probes (Finding 228) improve the combined Euclid DR1 significance:

\begin{center}
\begin{tabular}{lcc}
\toprule
\textbf{Probe Combination} & \textbf{Without new probes} & \textbf{With new probes} \\
\midrule
Euclid alone & $4.5\sigma$ & $\sim 5.0\sigma$ \\
Euclid + DESI Y3 & $5.1\sigma$ & $\sim 5.5\sigma$ \\
Euclid + DESI Y3 + Planck & $5.6\sigma$ & $\sim 6.0\sigma$ \\
\bottomrule
\end{tabular}
\end{center}

The improvement is modest ($\sim 0.4$--$0.5\sigma$) because the new probes are partially correlated with existing ones. But pushing from $4.5\sigma$ to $5.0\sigma$ for Euclid alone is strategically significant: it means $5\sigma$ detection may not require DESI Y3.

\section{Phase 2 Preparation Status}

\begin{center}
\begin{tabular}{lcl}
\toprule
\textbf{Paper} & \textbf{Readiness} & \textbf{Target Submission} \\
\midrule
$E_G$ prediction (JCAP) & 100\% & i88 \\
DESI standalone (JCAP) & 100\% & i88 \\
No-screening paper (PRD) & 100\% & i88 \\
Fermion masses (PRD) & 100\% & i88 \\
\bottomrule
\end{tabular}
\end{center}

All four Phase 2 papers are ready for submission at i88. This will bring the total submitted to 10 (1 arXiv + 9 journal submissions).

\section{Community Response Monitoring}

The $C_\ell$ paper (on arXiv since i71) has accumulated:
\begin{itemize}[nosep]
\item 12 citations (mostly self-citations from programme papers)
\item 3 external citations (one in a Euclid forecast paper, two in MOND reviews)
\item 0 critical responses or rebuttals
\item 2 private email inquiries (one asking about data availability, one asking about the $\psi$-screen mechanism)
\end{itemize}

The absence of critical responses is cautiously positive. The Euclid forecast paper citing DFD is particularly encouraging: it means the Euclid theory working group is aware of DFD predictions.

\end{document}
